\sezione{9}{Stima MLE in due dimensioni: esempio svolto}{44–49}

Stesso principio della sezione 6 (massima verosimiglianza), ora con vettori e matrici.

\subsection{I dati}
Training set di $n=5$ pattern bidimensionali:
\[
\left\{\,[3,\,7]\T,\ [4.5,\,12]\T,\ [13,\,7.2]\T,\ [1,\,4]\T,\ [23,\,15.8]\T\,\right\}
\]

\subsection{Le formule MLE}
\begin{formulat}{Stime di massima verosimiglianza}
$\displaystyle \mmu=\frac1n\sum_{i=1}^{n}\x_i
\qquad\qquad
\SSigma=\frac1n\sum_{i=1}^{n}(\x_i-\mmu)(\x_i-\mmu)\T$
\end{formulat}
\begin{intuizione}
La media vettoriale è la media \emph{componente per componente}. La covarianza è una somma di matrici: ogni pattern contribuisce con il prodotto «colonna per riga» del suo scarto dalla media, una matrice $2\times2$ che dice in quale direzione quel punto si allontana dal centro. Elemento per elemento:
\[
\sigma^{ij}=\frac1n\sum_{k=1}^{n}(x_k^i-\mu^i)(x_k^j-\mu^j)
\]
\end{intuizione}

\subsection{Il vettore medio}
\begin{esempio}
\[
\mmu=\begin{bmatrix}\dfrac{3+4.5+13+1+23}{5}\\[8pt]\dfrac{7+12+7.2+4+15.8}{5}\end{bmatrix}=\begin{bmatrix}8.9\\9.2\end{bmatrix}
\]
\end{esempio}

\subsection{La matrice di covarianza}
\begin{esempio}[Prodotti «colonna per riga» degli scarti]
\begin{center}\small
\rowcolors{2}{chiaro}{white}
\begin{tabular}{ccc}
\rowcolor{navy}\thd{$\x_i$} & \thd{$\x_i-\mmu$} & \thd{$(\x_i-\mmu)(\x_i-\mmu)\T$}\\
$[3,\,7]$ & $[-5.9,\,-2.2]$ & $\begin{psmallmatrix}34.81&12.98\\12.98&4.84\end{psmallmatrix}$\\[3pt]
$[4.5,\,12]$ & $[-4.4,\,2.8]$ & $\begin{psmallmatrix}19.36&-12.32\\-12.32&7.84\end{psmallmatrix}$\\[3pt]
$[13,\,7.2]$ & $[4.1,\,-2.0]$ & $\begin{psmallmatrix}16.81&-8.20\\-8.20&4.00\end{psmallmatrix}$\\[3pt]
$[1,\,4]$ & $[-7.9,\,-5.2]$ & $\begin{psmallmatrix}62.41&41.08\\41.08&27.04\end{psmallmatrix}$\\[3pt]
$[23,\,15.8]$ & $[14.1,\,6.6]$ & $\begin{psmallmatrix}198.81&93.06\\93.06&43.56\end{psmallmatrix}$\\
\end{tabular}
\end{center}
Sommando le cinque matrici e dividendo per 5:
\[
\SSigma=\begin{bmatrix}\sigma^{11}&\sigma^{12}\\\sigma^{21}&\sigma^{22}\end{bmatrix}=\begin{bmatrix}66.44&25.32\\25.32&17.456\end{bmatrix}
\]
Per esempio $\sigma^{11}=(34.81+19.36+16.81+62.41+198.81)/5=332.2/5=66.44$.
\end{esempio}
\textbf{Lettura:} $\sigma^{12}>0$, quindi le due feature crescono insieme e l'ellisse è inclinata verso l'alto a destra. La varianza lungo $x^1$ (66.4) è molto maggiore di quella lungo $x^2$ (17.5): l'ellisse è allungata in orizzontale.

\subsection{Determinante e inversa}
\begin{esempio}
\[
|\SSigma|=66.44\cdot17.456-25.32^2=1159.78-641.10=518.674
\]
\[
\SSigma^{-1}=\frac{1}{518.674}\begin{bmatrix}17.456&-25.32\\-25.32&66.44\end{bmatrix}=\begin{bmatrix}0.0337&-0.0488\\-0.0488&0.1281\end{bmatrix}
\]
È la formula $2\times2$ della sezione 7: si scambiano gli elementi della diagonale, si cambia segno agli altri, si divide per il determinante.
\end{esempio}
\begin{tecnica}[In più: autovalori di questa $\SSigma$]
$\lambda_1\approx77.18,\ \lambda_2\approx6.72$. Verifiche: $\lambda_1+\lambda_2=83.90=66.44+17.456$ (traccia) e $\lambda_1\lambda_2\approx518.67=|\SSigma|$. L'asse principale punta a circa $23^\circ$ (direzione $\approx(0.92,\,0.39)$): l'ellisse è lunga $\sqrt{77.18}\approx8.8$ e larga $\sqrt{6.72}\approx2.6$ (in unità di $r=1$).
\end{tecnica}

\fig[.98\linewidth]{stima2d.pdf}
\didascalia{La gaussiana stimata vista dall'alto (curve di livello, con gli autovettori scalati per $2\sqrt{\lambda}$) e di lato (la campana bidimensionale). Corrisponde alla slide 49.}

\begin{attenzione}
\begin{itemize}
  \item Come in 1D, la MLE divide per $n$; con $n-1$ si otterrebbe $\begin{psmallmatrix}83.05&31.65\\31.65&21.82\end{psmallmatrix}$. Con soli 5 punti la differenza è grande (25\%).
  \item Con 5 punti in 2D la stima è molto incerta: il punto $[23,\,15.8]$ da solo pesa il 60\% di $\sigma^{11}$. Servono $n>d$ per l'invertibilità e molti più punti per una stima affidabile.
\end{itemize}
\end{attenzione}
\begin{mlbox}[Calcolo incrementale]
Vale l'identità $\SSigma=\frac1n\sum_i\x_i\x_i\T-\mmu\mmu\T$: basta tenere le somme $\sum_i\x_i$ e $\sum_i\x_i\x_i\T$ e aggiornarle a ogni nuovo esempio, senza conservare i dati. È la proprietà «incrementale» della sezione 3.
\end{mlbox}

\begin{appunti}{— Sezione 9}
\pag{44} 5 pattern 2D; stima MLE di vettore medio e covarianza.
\pag{45} $\mmu=\frac1n\sum_i\x_i=[8.9,\ 9.2]\T$ (media componente per componente).
\pag{46} $\SSigma=\frac1n\sum_i(\x_i-\mmu)(\x_i-\mmu)\T$: somma di matrici colonna $\times$ riga; $\sigma^{ij}=\frac1n\sum_k(x_k^i-\mu^i)(x_k^j-\mu^j)$.
\pag{47} $\SSigma=\begin{psmallmatrix}66.44&25.32\\25.32&17.456\end{psmallmatrix}$; $|\SSigma|=518.674$; $\SSigma^{-1}=\begin{psmallmatrix}0.0337&-0.0488\\-0.0488&0.1281\end{psmallmatrix}$. $\sigma^{12}>0$: ellisse inclinata.
\pag{48} Inversa $2\times2$: $\frac1{ad-bc}\begin{psmallmatrix}d&-b\\-c&a\end{psmallmatrix}$.
\pag{49} Autovalori $\approx77.2,\ 6.7$; asse principale a $\approx23^\circ$. Con $n-1$: stime più grandi del 25\%. $\SSigma=\frac1n\sum\x\x\T-\mmu\mmu\T$ (incrementale).
\end{appunti}
